\chapter{Catheter Problems}
\index{Catheter Problems}

\section*{Definition and Overview}
A urinary catheter is a thin tube inserted into the bladder through the urethra, or occasionally directly through the abdominal wall, to drain urine. Catheters are used for a range of reasons, including bladder obstruction, surgery, monitoring of urine output, or incontinence that cannot be managed otherwise. Most catheters are temporary, but some people live with long-term or indwelling catheters for months or years.

Catheter problems are common and include infection, blockage, leaking, discomfort, and accidental removal. Managing these problems involves a combination of good hygiene, regular catheter care, prompt treatment of infection, and knowing when to seek urgent help. A catheter that is blocked or causing fever is a medical issue that requires prompt attention.

\section*{Epidemiology}
Catheter problems are among the most common complications of hospital and long-term care. Catheter-associated urinary tract infection (CAUTI) is the leading cause of healthcare-associated infection worldwide, and the risk of infection increases the longer a catheter stays in place. Around half of people with a long-term catheter develop a urinary infection within weeks, and most develop recurrent infections over time.

Blockage is also very common in long-term catheter users, with some people experiencing repeated blockages that require catheter changes or flushing. Leaking around the catheter (bypassing) affects many users and is often a sign of blockage, spasm, or a catheter that is too large. These problems cause significant discomfort, hospital visits, and health costs.

\section*{Aetiology and Pathophysiology}
Catheter problems arise from several mechanisms:

\begin{itemize}
    \item \textbf{Infection:} bacteria travel along the inside or outside of the catheter into the bladder, forming biofilms that are hard to clear with antibiotics
    \item \textbf{Blockage:} sediment, blood clots, or encrustation from mineral deposits can obstruct the catheter lumen or the drainage bag
    \item \textbf{Bladder spasm:} the catheter irritates the bladder wall, causing urgency, leaking, or cramping
    \item \textbf{Bypassing:} urine leaks around the catheter, often because of blockage, spasm, or an incorrect catheter size
    \item \textbf{Mechanical problems:} kinking of the tube, a blocked or displaced drainage bag, or pulling on the catheter
    \item \textbf{Trauma:} the catheter can irritate, inflame, or injure the urethra or bladder lining
\end{itemize}

The presence of a foreign body in the bladder impairs the bladder's natural defences, and bacteria can multiply freely on the catheter surface, which is why infections recur so easily.

\section*{Clinical Features}
\begin{itemize}
    \item Fever, chills, or feeling unwell, suggesting a urinary infection
    \item Cloudy, smelly, or blood-stained urine
    \item Burning, stinging, or discomfort in the bladder or urethra
    \item Pain or pressure in the lower abdomen or back
    \item Leaking of urine around the catheter (bypassing)
    \item Little or no urine draining into the bag, suggesting blockage
    \item Visible sediment, grit, or particles in the catheter or bag
    \item Urgency, cramping, or bladder spasms
    \item A catheter that has fallen out or been pulled out accidentally
\end{itemize}

\section*{Differential Diagnosis}
\begin{itemize}
    \item \textbf{Catheter-associated urinary tract infection} vs.\ \textbf{colonisation:} bacteria in the urine without symptoms usually need no treatment
    \item \textbf{Blockage} vs.\ \textbf{mechanical kink or drainage bag problem} causing low urine output
    \item \textbf{Bladder spasm} vs.\ \textbf{infection} as the cause of urgency and pain
    \item \textbf{Bypassing} due to blockage, spasm, or an over-sized catheter
    \item \textbf{Pyelonephritis or sepsis:} an infection that has spread to the kidney or bloodstream, causing fever, rigors, and loin pain
    \item \textbf{Trauma:} bleeding or urethral injury from the catheter itself
\end{itemize}

\section*{When to Seek Medical Care}
Contact a doctor or nurse, or the team that manages your catheter, if you have fever or chills, severe pain, no urine draining for several hours, blood in the urine, or a catheter that has fallen out and cannot be replaced. Simple bypassing, sediment, or mild irritation can often be managed with advice, but it is safest to have any new problem reviewed.

\textbf{Contact your local emergency services for:} high fever with rigors, confusion or altered consciousness, severe abdominal or loin pain, inability to pass any urine with a blocked catheter, or signs of septic shock such as rapid breathing, fast heartbeat, or feeling faint.

\section*{Diagnosis and Investigations}
\begin{itemize}
    \item \textbf{History:} how long the catheter has been in place, the exact nature of the problem, and any previous episodes
    \item \textbf{Examination:} checking the catheter position, the drainage bag, the abdomen for distension or tenderness, and vital signs for signs of infection
    \item \textbf{Urine dipstick and culture:} identifying bacteria and confirming whether infection is present; a urine sample is taken from the catheter port, not the bag
    \item \textbf{Blood tests:} full blood count, inflammatory markers, and kidney function in suspected infection or obstruction
    \item \textbf{Ultrasound or scan:} assessment of the bladder and kidneys where blockage or kidney problems are suspected
    \item \textbf{Flushing:} gently flushing the catheter with sterile fluid can confirm and sometimes clear a blockage
\end{itemize}

\section*{Management}
\begin{itemize}
    \item \textbf{Infection:} antibiotics selected on the basis of urine culture, together with ensuring good fluid intake; catheter-associated infection often improves when the catheter is changed or removed
    \item \textbf{Blockage:} catheter flushing, increasing fluid intake, and possibly changing to a catheter or regimen better suited to the person; frequent blockers may benefit from a different catheter material or a regular flushing schedule
    \item \textbf{Bladder spasm:} antispasmodic medicines can reduce urgency, cramping, and leaking
    \item \textbf{Bypassing:} usually managed by treating the underlying cause, such as blockage or spasm, or by reviewing catheter size
    \item \textbf{Catheter removal or change:} short-term catheters should be removed as soon as they are no longer needed, and changed on a schedule for long-term users
    \item \textbf{Drainage bag care:} emptying, cleaning, and positioning the bag correctly to prevent backflow
    \item \textbf{Alternative approaches:} intermittent self-catheterisation or a suprapubic catheter may suit some people better than a long-term indwelling urethral catheter
\end{itemize}

\section*{Complications}
\begin{itemize}
    \item Catheter-associated urinary tract infection, including kidney infection and bloodstream infection (sepsis)
    \item Chronic inflammation and bleeding of the bladder lining (cystitis)
    \item Urethral injury, stricture, or erosion from long-term catheter use
    \item Bladder stones forming around the catheter tip
    \item Repeated blockages leading to urine retention and acute kidney injury in severe cases
    \item Discomfort, leaking, and loss of dignity and independence
    \item The longer a catheter remains in place, the higher the cumulative risk of these complications
\end{itemize}

\section*{Prognosis and Outcomes}
Most acute catheter problems resolve promptly with the right treatment: infections settle with antibiotics, and blockages clear with flushing or a catheter change. For people who need long-term catheters, many problems can be anticipated and managed with good routines, and most people continue to manage well with planned catheter changes and simple hygiene measures. The outlook is best when catheters are used only for as long as needed, and when infections and blockages are treated early rather than left to become severe.

\section*{Prevention and Self-Care}
\begin{itemize}
    \item Wash hands before and after handling the catheter or drainage bag
    \item Keep the catheter area clean and dry, and clean around the entry site daily
    \item Keep the drainage bag below the level of the bladder and never lift it above the waist
    \item Empty the bag regularly and clean it as advised by the health team
    \item Drink plenty of fluids unless advised otherwise, to keep urine dilute and flowing
    \item Avoid kinks, pulls, and tension on the catheter tubing
    \item Use a catheter securement device or leg strap to hold the catheter in place
    \item Remove or change the catheter as soon as recommended, and never keep a short-term catheter longer than needed
    \item Seek advice early for fever, pain, or reduced urine output rather than waiting
\end{itemize}

\section*{Sources and Further Reading}
The following organisations provide reliable information on urinary catheters and catheter care.

\begin{itemize}
    \item National Health Service. \textit{NHS conditions guide on urinary catheters and their care.} https://www.nhs.uk/conditions/
    \item Mayo Clinic. \textit{Urinary catheter: care and troubleshooting information.} https://www.mayoclinic.org/
    \item Centers for Disease Control and Prevention. \textit{Catheter-associated urinary tract infection (CAUTI) prevention guidance.} https://www.cdc.gov/
    \item World Health Organization. \textit{Healthcare-associated infection prevention and control guidance.} https://www.who.int/
    \item MedlinePlus. \textit{Consumer health information on urinary catheter care.} https://medlineplus.gov/
    \item MSD Manual. \textit{Urinary catheterisation and its complications -- professional overview.} https://www.msdmanuals.com/
\end{itemize}
